Considera la funció
\[
f(x)=\frac{x}{x^2+1} .
\]

\begin{apartats}

\apartat[1,25]{1}
Estudia la monotonia de $f$ i troba'n els extrems relatius.

\begin{solucio}
$f'(x)=\dfrac{\left(x^2+1\right)-x\cdot2x}{\left(x^2+1\right)^2}=\dfrac{1-x^2}{\left(x^2+1\right)^2}$.\\
El denominador és sempre positiu, així que el signe és el de $1-x^2$, que s'anul·la en
$x=\pm1$.\\
$f$ \textbf{decreix} a $(-\infty,-1)$, \textbf{creix} a $(-1,1)$ i \textbf{decreix} a
$(1,+\infty)$.\\
\textbf{Mínim relatiu} $\left(-1,-\tfrac12\right)$ i \textbf{màxim relatiu}
$\left(1,\tfrac12\right)$.
\end{solucio}

\begin{nomesllarg}
\apartat{0,75}
Quins són els valors màxim i mínim de $f$ a l'interval $[-3,3]$?

\begin{solucio}
Cal comparar els extrems relatius de dins amb els valors dels extrems de l'interval:
$f(-3)=-\dfrac{3}{10}$, $f(-1)=-\dfrac12$, $f(1)=\dfrac12$ i $f(3)=\dfrac{3}{10}$.\\
El valor \textbf{màxim} és $\dfrac12$, en $x=1$, i el \textbf{mínim} és $-\dfrac12$, en
$x=-1$.
\end{solucio}
\end{nomesllarg}

\apartat[1,25]{0,75}
Demostra que la funció $h(x)=-x^3-x+4$ és decreixent a tot $\mathbb{R}$.

\begin{solucio}
$h'(x)=-3x^2-1$. Com que $-3x^2\le0$, es compleix $h'(x)\le-1<0$ per a tot $x$.\\
La derivada és sempre negativa, i per tant $h$ \textbf{decreix a tot $\mathbb{R}$}.
\end{solucio}

\end{apartats}
